\begin{abstract}
KV cache compression methods like H2O eviction and quantization enable memory-efficient long-context inference, yet they apply uniform strategies across all tasks. We ask whether task structure should inform compression selection. Through gap statistic analysis of 21 LongBench tasks across 6 compression configurations, we discover that tasks cluster into three statistically distinct response groups ($k^* = 3$, gap exceeding standard error), with interpretable structure: multi-document QA and code tasks show high eviction sensitivity, while summarization tasks tolerate aggressive compression. We further demonstrate that first-100-token attention entropy discriminates task domains with large effect size (F = 38.05, $p = 2.92 \times 10^{-26}$, $\eta^2 = 0.522$), validating entropy as a potential routing signal. These findings establish the empirical foundation for task-conditioned KV cache compression: task-dependent structure exists and is detectable via attention features. While the complete routing system remains future work, our results challenge the one-size-fits-all assumption underlying current methods and motivate task-aware compression selection.
\end{abstract}
